\clearpage
\section{03 关系模型与完整性约束：中文学习笔记}\label{03_relational_model_and_integrity_zh}

CSC5041 · 第一次课堂测验复习 · 2026-10-19

\subsection{阅读说明与学习目标}\label{03_relational_model_and_integrity_zh-ux9605ux8bfbux8bf4ux660eux4e0eux5b66ux4e60ux76eeux6807}

本篇接续 \hyperref[02_er_modelling_zh]{02 ER 概念建模}，对应教学计划第 1 页 Week 3 的 ``The Relational Data Model'' 与 ``integrity constraints''。主要依据教师 \href{../../../slides/Lecture2A.pptx}{Lecture2A.pptx} 第 4--32 页；下文未注明文件的页码均按该 PPT 的幻灯片顺序计数。

学完后，应能读懂关系模式和实例，解释关系的集合性质，找出各类键及外键引用方向，区分键约束、实体完整性与参照完整性，并判断插入、删除、修改是否破坏约束。第 33 页开始的 ER 映射留到 04；关系代数留到 05；SQL 实现细节留到 06。

文中的学生选课数据是\textbf{自编教学示例}，不是真实学生资料、运行结果或考试题。对课件简化表述的辨析会明确标出。

\subsection{1. 从 ER 模型走到关系模型}\label{03_relational_model_and_integrity_zh-ux4ece-er-ux6a21ux578bux8d70ux5230ux5173ux7cfbux6a21ux578b}

来源：第 4--7、32 页。

02 用实体、属性和关系表达应用含义；关系模型用一组相互关联的\textbf{关系 relation} 组织数据。这里的 relation 可直观理解为表，\textbf{不是 ER 图中连接实体的 relationship}。

\begin{longtable}[]{@{}lll@{}}
\toprule\noalign{}
数学术语 & 常见数据术语 & 含义 \\
\midrule\noalign{}
\endhead
\bottomrule\noalign{}
\endlastfoot
Relation 关系 & Table 表 & 一组具有相同属性结构的元组 \\
Tuple 元组 & Row / record 行、记录 & 某组属性的一组取值 \\
Attribute 属性 & Column / field 列、字段 & 被命名的性质 \\
Domain 域 & Allowed values 允许值集合 & 该属性允许取得哪些值 \\
\end{longtable}

例如 \texttt{Student(\allowbreak{}sid,\allowbreak{}\ name,\allowbreak{}\ email)\allowbreak{}} 是一个关系模式；\texttt{(\allowbreak{}S01,\allowbreak{}\ Amy,\allowbreak{}\ amy@example.\allowbreak{}test)\allowbreak{}} 是按该属性顺序书写的一个元组。Student 与 Course 之间的选课联系也可以由一个单独的关系 Enrolment 表示，不是靠把两张表放在相邻位置表示。

\subsection{2. 模式、实例、度与基数}\label{03_relational_model_and_integrity_zh-ux6a21ux5f0fux5b9eux4f8bux5ea6ux4e0eux57faux6570}

来源：第 8、32 页；以下形式化记法与 degree / cardinality 术语是解释性补充。

\subsubsection{单个关系}\label{03_relational_model_and_integrity_zh-ux5355ux4e2aux5173ux7cfb}

\begin{StudyCode}
关系模式：Student(sid, name, email)

一个关系实例：
sid   name   email
S01   Amy    amy@example.test
S02   Ben    ben@example.test
\end{StudyCode}

关系模式规定关系名、属性及其域；关系实例是某一时刻该结构下的实际元组集合。常写作 \texttt{R(\allowbreak{}A1,\allowbreak{}\ .\allowbreak{}.\allowbreak{}.\allowbreak{},\allowbreak{}\ An)\allowbreak{}} 与 \texttt{r(\allowbreak{}R)\allowbreak{}}，或在上下文清楚时简写为 \texttt{r}。

不考虑 NULL 的经典定义中，每个元组从各属性域各取一个值，因此 \(r\,\subseteq \,\mathsf{dom}(\mathsf{A1})\,\times \,...\,\times \,\mathsf{dom}(\mathsf{An})\)。符号只说明``每列从哪里取值、实例是哪些组合的集合''，不要求把所有可能组合都存进去。

\begin{longtable}[]{@{}lll@{}}
\toprule\noalign{}
术语 & 数什么？ & 上例 \\
\midrule\noalign{}
\endhead
\bottomrule\noalign{}
\endlastfoot
度 Degree / arity & 属性数 & 3 \\
基数 Cardinality & 当前元组数 & 2 \\
\end{longtable}

空实例仍有模式：Student 暂时没有记录时，度仍为 3，基数为 0。这里的 cardinality 是\textbf{当前关系有多少行}；02 的 ER cardinality ratio 是\textbf{业务规则允许关联多少个对方实体}，不要混用。

\subsubsection{整个数据库}\label{03_relational_model_and_integrity_zh-ux6574ux4e2aux6570ux636eux5e93}

第 32 页将关系数据库模式定义为\textbf{一组关系模式及一组完整性约束}。因此仅列 Student、Course、Enrolment 的列名，还没有完整写出数据库模式；也要说明主键、外键等规则。

数据库状态包含各关系的实例；全部相关约束都满足时，才是合法状态。以下讨论的``非法状态''用来说明操作为何应被拒绝，并不表示系统应该保存它。

\subsection{3. 一个关系的基本性质}\label{03_relational_model_and_integrity_zh-ux4e00ux4e2aux5173ux7cfbux7684ux57faux672cux6027ux8d28}

来源：第 8--10 页；SQL 对照另见 SQLpart1 第 31--32 页。

\begin{longtable}[]{@{}
  >{\raggedright\arraybackslash}p{(\linewidth - 2\tabcolsep) * \real{0.5000}}
  >{\raggedright\arraybackslash}p{(\linewidth - 2\tabcolsep) * \real{0.5000}}@{}}
\toprule\noalign{}
\begin{minipage}[b]{\linewidth}\raggedright
性质
\end{minipage} & \begin{minipage}[b]{\linewidth}\raggedright
正确理解
\end{minipage} \\
\midrule\noalign{}
\endhead
\bottomrule\noalign{}
\endlastfoot
元组无顺序 & 交换两行的显示位置，不会得到另一个关系 \\
元组不能完全重复 & 经典关系是集合，相同元组不作为两个不同元素出现 \\
每个属性位置取一个原子值 & 本课的平面关系模型不直接保留 ER 复合或多值属性 \\
属性按名称区分 & 同一关系的属性名需要能区分各列；不同关系可有同名属性 \\
重排显示列不改变含义 & 必须把列名和各自的值一起移动，不能只随意交换元组中的值 \\
\end{longtable}

\subsubsection{列无语义顺序，不等于元组值可以乱放}\label{03_relational_model_and_integrity_zh-ux5217ux65e0ux8bedux4e49ux987aux5e8fux4e0dux7b49ux4e8eux5143ux7ec4ux503cux53efux4ee5ux4e71ux653e}

第 10 页将 PLAYER 的列连同列名与值一起重排，得到相同关系。用命名属性理解时，重要的是``哪个值属于哪个属性''。

但在 \texttt{Student(\allowbreak{}sid,\allowbreak{}\ name)\allowbreak{}} 下，\texttt{(\allowbreak{}S01,\allowbreak{}\ Amy)\allowbreak{}} 与 \texttt{(\allowbreak{}Amy,\allowbreak{}\ S01)\allowbreak{}} 不一样：按位置书写元组时，位置必须与所声明的属性顺序对应。重排展示与交换属性值是两件事。

\subsubsection{原子值与 ER 属性}\label{03_relational_model_and_integrity_zh-ux539fux5b50ux503cux4e0e-er-ux5c5eux6027}

02 中的 \texttt{address(\allowbreak{}city,\allowbreak{}\ street,\allowbreak{}\ postcode)\allowbreak{}} 是复合属性，\texttt{\{phoneNumber\}} 是多值属性。第 8 页明确排除在本课平面关系中直接使用这两种属性结构。原子性取决于当前模型所采用的域与操作，不是指一个字符串在字符层面无法拆开。具体怎样转换留到 04。

\subsubsection{经典关系与 SQL 表}\label{03_relational_model_and_integrity_zh-ux7ecfux5178ux5173ux7cfbux4e0e-sql-ux8868}

经典关系遵循集合语义；SQLpart1 第 31--32 页另指出 SQL 可以保留重复，查询结果可用 DISTINCT 去重。因此不能把``关系无重复''直接当作``任何 SQL 表都自动禁止重复''的结论。SQL 语法与具体约束声明留到 06。

\subsection{4. 超键、候选键与主键}\label{03_relational_model_and_integrity_zh-ux8d85ux952eux5019ux9009ux952eux4e0eux4e3bux952e}

来源：第 12--15 页；与 02 第 4 节衔接。

\begin{longtable}[]{@{}
  >{\raggedright\arraybackslash}p{(\linewidth - 4\tabcolsep) * \real{0.3333}}
  >{\raggedright\arraybackslash}p{(\linewidth - 4\tabcolsep) * \real{0.3333}}
  >{\raggedright\arraybackslash}p{(\linewidth - 4\tabcolsep) * \real{0.3333}}@{}}
\toprule\noalign{}
\begin{minipage}[b]{\linewidth}\raggedright
键
\end{minipage} & \begin{minipage}[b]{\linewidth}\raggedright
条件
\end{minipage} & \begin{minipage}[b]{\linewidth}\raggedright
自编例子：假定 sid 与 email 分别唯一且必填
\end{minipage} \\
\midrule\noalign{}
\endhead
\bottomrule\noalign{}
\endlastfoot
超键 Superkey & 能唯一确定一个元组 & \texttt{\{sid\}}、\texttt{\{email\}}、\texttt{\{sid,\ name\}} \\
候选键 Candidate key & 是超键，且没有多余属性 & \texttt{\{sid\}}、\texttt{\{email\}} \\
主键 Primary key & 被选定的一个候选键 & 选择 \texttt{\{sid\}} \\
复合键 Composite key & 由多个属性组成 & Course 的 \texttt{\{dept,\ courseNo\}}，见第 8 节 \\
\end{longtable}

最小性检查的是\textbf{删掉属性后还能否保证唯一}，不是``在全体候选键中挑属性数最少的''。一张表可以有多个候选键，选定一个主键也不意味着其他候选键失去唯一性要求。

第 13 页以``没有两人同名''为\textbf{显式假设}，才把 PLAYER 的 Name 当候选键；不能从示例的三个人恰好不同名推导现实中姓名总是唯一。

\subsubsection{为什么无重复元组还需要键约束？}\label{03_relational_model_and_integrity_zh-ux4e3aux4ec0ux4e48ux65e0ux91cdux590dux5143ux7ec4ux8fd8ux9700ux8981ux952eux7ea6ux675f}

\begin{StudyCode}
(S01, Amy, amy@example.test)
(S01, Ada, ada@example.test)
\end{StudyCode}

这两个完整元组不同，集合性质并不排除它们同时出现。但是它们具有相同 sid，若 sid 是主键，就违反键约束。\textbf{无完全重复行}与\textbf{指定键不能重复}是不同要求。

\subsubsection{\texorpdfstring{第 15 页的 \texttt{2\^{}n\ −\ 1} 怎么理解？}{第 15 页的 2\^{}n − 1 怎么理解？}}\label{03_relational_model_and_integrity_zh-ux7b2c-15-ux9875ux7684-2n-1-ux600eux4e48ux7406ux89e3}

该页列出所有非空属性子集。对 n 个属性，非空子集确实有 \texttt{2\^{}n\ −\ 1} 个，\textbf{但不是每个子集都能唯一识别元组}，所以不能将它当作任意关系的超键数公式。

\textbf{自编反例}：R(A,B,C) 只有候选键 \texttt{\{A\}}，B、C 及 \texttt{\{B,C\}} 都不保证唯一。超键只有 \texttt{\{A\}}、\texttt{\{A,B\}}、\texttt{\{A,C\}}、\texttt{\{A,B,C\}}，共 4 个，不是 7 个。只有在每个单属性本身都能作键等额外条件下，所有非空子集才全是超键。

\Needspace{11cm}
\subsection{5. 外键：方向、重复值与复合引用}\label{03_relational_model_and_integrity_zh-ux5916ux952eux65b9ux5411ux91cdux590dux503cux4e0eux590dux5408ux5f15ux7528}

\StudyVisual{integrity}{zh}

来源：第 16--19 页。

外键 Foreign key 是引用方关系中的一个或一组属性，用其取值引用被引用关系的键。课件本节采用引用\textbf{主键}的定义；本篇按这一范围讲解，不把它扩展成所有 SQL 系统的完整规则。

\begin{StudyCode}
引用方 referencing relation                 被引用方 referenced relation
Enrolment.sid                  ─────────→   Student.sid
Enrolment.(dept, courseNo)      ─────────→   Course.(dept, courseNo)
\end{StudyCode}

核对外键时，要写清楚\textbf{引用方属性组 → 目标关系的键属性组}。不要求两边属性名相同，但对应域应兼容；仅仅都是字符串或列名相同，不能自动成立外键。

\subsubsection{外键自身不要求唯一}\label{03_relational_model_and_integrity_zh-ux5916ux952eux81eaux8eabux4e0dux8981ux6c42ux552fux4e00}

多条选课记录可以有相同 sid，表示同一名学生选了多门课。Student.sid 必须唯一，Enrolment.sid 作为外键则可以重复。外键也可以同时是主键的一部分，这两个身份不冲突。

\subsubsection{复合外键必须作为一个组合匹配}\label{03_relational_model_and_integrity_zh-ux590dux5408ux5916ux952eux5fc5ux987bux4f5cux4e3aux4e00ux4e2aux7ec4ux5408ux5339ux914d}

若 Course 中只有 \texttt{(\allowbreak{}CS,\allowbreak{}101)\allowbreak{}} 和 \texttt{(\allowbreak{}Math,\allowbreak{}201)\allowbreak{}}，那么 \texttt{(\allowbreak{}CS,\allowbreak{}201)\allowbreak{}} \textbf{不能}引用它们。虽然 CS 和 201 分别出现过，但它们没有在\textbf{同一个目标元组的对应键位置}共同出现。

课件第 19 页中的 \texttt{(\allowbreak{}Department,\allowbreak{}\ Degree)\allowbreak{}} 也是\textbf{一组复合外键}，不是 Department 和 Degree 各自单独引用 COURSE 的两个外键。遗漏任一部分，可能无法确定具体课程。

\subsection{6. NULL 与外键的边界}\label{03_relational_model_and_integrity_zh-null-ux4e0eux5916ux952eux7684ux8fb9ux754c}

来源：第 18、20、22、30 页；NULL 的未知/不适用含义已在 02 解释。

\begin{longtable}[]{@{}ll@{}}
\toprule\noalign{}
所在位置 & 本篇采用的规则 \\
\midrule\noalign{}
\endhead
\bottomrule\noalign{}
\endlastfoot
主键的任何组成属性 & 不能为 NULL，即使其他组成属性已有值 \\
非主键属性 & 是否允许 NULL 取决于该属性的必填约束 \\
可空外键 & 只有在其他约束也允许时，才可用 NULL 表示暂不引用 \\
同时属于主键的外键属性 & 不能因为它也是外键就放宽主键非空规则 \\
\end{longtable}

\textbf{课件第 18、22 页的复合外键规则是：要么全部 NULL，要么全部非 NULL 且完整匹配目标键。} 本篇按这个明确给出的课堂规则判断，不将其泛化为所有 SQL 实现的默认行为。

\textbf{自编小例子}：给定 \texttt{Plan(\allowbreak{}planID,\allowbreak{}\ dept,\allowbreak{}\ courseNo)\allowbreak{}}，主键为 planID，\texttt{(\allowbreak{}dept,\allowbreak{}courseNo)\allowbreak{}} 是可空且遵守上述规则的外键，引用第 5 节的 Course。

\begin{longtable}[]{@{}ll@{}}
\toprule\noalign{}
外键值 & 判定 \\
\midrule\noalign{}
\endhead
\bottomrule\noalign{}
\endlastfoot
\texttt{(\allowbreak{}CS,\allowbreak{}101)\allowbreak{}} & 合法，完整匹配 \\
\texttt{(\allowbreak{}NULL,\allowbreak{}NULL)\allowbreak{}} & 合法，假定两列都允许缺失且没有其他必填要求 \\
\texttt{(\allowbreak{}CS,\allowbreak{}NULL)\allowbreak{}} & 不合法，不符合本课``全空或全非空''的规则 \\
\texttt{(\allowbreak{}CS,\allowbreak{}201)\allowbreak{}} & 不合法，不存在这个组合 \\
\end{longtable}

\textbf{``参照完整性允许空引用''不等于``业务上允许没有关联''。} 若要求每条记录必须引用目标，还需要禁止外键为空。02 的全参与约束不能只靠``有一个可空外键''完整表达。

\subsection{7. 三类完整性约束检查什么？}\label{03_relational_model_and_integrity_zh-ux4e09ux7c7bux5b8cux6574ux6027ux7ea6ux675fux68c0ux67e5ux4ec0ux4e48}

来源：第 20--22、32 页。

\begin{longtable}[]{@{}
  >{\raggedright\arraybackslash}p{(\linewidth - 4\tabcolsep) * \real{0.3333}}
  >{\raggedright\arraybackslash}p{(\linewidth - 4\tabcolsep) * \real{0.3333}}
  >{\raggedright\arraybackslash}p{(\linewidth - 4\tabcolsep) * \real{0.3333}}@{}}
\toprule\noalign{}
\begin{minipage}[b]{\linewidth}\raggedright
约束
\end{minipage} & \begin{minipage}[b]{\linewidth}\raggedright
要防止的问题
\end{minipage} & \begin{minipage}[b]{\linewidth}\raggedright
典型违规
\end{minipage} \\
\midrule\noalign{}
\endhead
\bottomrule\noalign{}
\endlastfoot
键约束 Key constraint & 不同元组共享同一个候选键值 & 两名学生使用相同 sid；或 sid 不同但唯一 email 重复 \\
实体完整性 Entity integrity & 主键不能完整识别一个元组 & 复合主键中某一属性为 NULL \\
参照完整性 Referential integrity & 外键引用不存在的目标键 & 选课记录引用数据库中没有的学生 \\
\end{longtable}

``主键唯一''属于键约束，``主键非空''属于实体完整性。它们共同支持标识，但答题时应分别说明。一个操作也可能同时违反多条规则；不要找到第一项就认定没有其他问题。

\subsubsection{域约束和业务约束【解释性补充】}\label{03_relational_model_and_integrity_zh-ux57dfux7ea6ux675fux548cux4e1aux52a1ux7ea6ux675fux89e3ux91caux6027ux8865ux5145}

第 6、8 页的域规定值的范围。例如 grade 必须取规定等级或允许缺失，整数编号不能随意填自由文本。其他业务规则还可能要求费用非负、开始日期不晚于结束日期。

本篇三类约束是课件第 20 页重点列出的约束，\textbf{不是数据库全部可能约束的清单}。仅主键和外键全部合法，也不保证业务状态一定正确。

\subsection{8. 贯穿例子：先明确模式、约束与初始状态}\label{03_relational_model_and_integrity_zh-ux8d2fux7a7fux4f8bux5b50ux5148ux660eux786eux6a21ux5f0fux7ea6ux675fux4e0eux521dux59cbux72b6ux6001}

以下全部为自编教学例子。模式是已给定输入，不在这里演示 ER 映射。

\subsubsection{模式与约束}\label{03_relational_model_and_integrity_zh-ux6a21ux5f0fux4e0eux7ea6ux675f}

\begin{longtable}[]{@{}
  >{\raggedright\arraybackslash}p{(\linewidth - 6\tabcolsep) * \real{0.2500}}
  >{\raggedright\arraybackslash}p{(\linewidth - 6\tabcolsep) * \real{0.2500}}
  >{\raggedright\arraybackslash}p{(\linewidth - 6\tabcolsep) * \real{0.2500}}
  >{\raggedright\arraybackslash}p{(\linewidth - 6\tabcolsep) * \real{0.2500}}@{}}
\toprule\noalign{}
\begin{minipage}[b]{\linewidth}\raggedright
关系模式
\end{minipage} & \begin{minipage}[b]{\linewidth}\raggedright
主键 PK
\end{minipage} & \begin{minipage}[b]{\linewidth}\raggedright
其他候选键
\end{minipage} & \begin{minipage}[b]{\linewidth}\raggedright
外键 FK
\end{minipage} \\
\midrule\noalign{}
\endhead
\bottomrule\noalign{}
\endlastfoot
Student(sid, name, email) & \texttt{\{sid\}} & \texttt{\{email\}} & 无 \\
Course(dept, courseNo, title) & \texttt{\{dept,\ courseNo\}} & 本例未指定其他候选键 & 无 \\
Enrolment(sid, dept, courseNo, term, grade) & \texttt{\{sid,\ dept,\ courseNo,\ term\}} & 本例未指定其他候选键 & \texttt{\{sid\}\ →\ Student.sid}；\texttt{\{dept,courseNo\}\ →\ Course.\allowbreak{}(dept,\allowbreak{}courseNo)} \\
\end{longtable}

本例规定：Student 的三个属性都必填；课程三个属性都必填；sid、email、dept、title、name、term 是字符串，courseNo 是正整数；term 允许 \texttt{2026F}、\texttt{2027S}；grade 允许 A、B、C、D、F 或 NULL。NULL 表示尚未获得成绩。所有主键组成属性均非空。

一名学生可以选多门课，也可在不同 term 再选同一门课；\textbf{同一学生、同一课程、同一 term 最多一条选课记录}。这正是 Enrolment 四属性主键表达的规则。dept 或 courseNo 单独都不能唯一识别 Course。

\subsubsection{初始状态}\label{03_relational_model_and_integrity_zh-ux521dux59cbux72b6ux6001}

Student：

\begin{longtable}[]{@{}lll@{}}
\toprule\noalign{}
sid & name & email \\
\midrule\noalign{}
\endhead
\bottomrule\noalign{}
\endlastfoot
S01 & Amy & amy@example.test \\
S02 & Ben & ben@example.test \\
\end{longtable}

Course：

\begin{longtable}[]{@{}lll@{}}
\toprule\noalign{}
dept & courseNo & title \\
\midrule\noalign{}
\endhead
\bottomrule\noalign{}
\endlastfoot
CS & 101 & Database Basics \\
Math & 201 & Linear Algebra \\
\end{longtable}

Enrolment：

\begin{longtable}[]{@{}lllll@{}}
\toprule\noalign{}
sid & dept & courseNo & term & grade \\
\midrule\noalign{}
\endhead
\bottomrule\noalign{}
\endlastfoot
S01 & CS & 101 & 2026F & A \\
S02 & CS & 101 & 2026F & NULL \\
\end{longtable}

这个状态满足本例列出的约束：Student 两个候选键都唯一；所有主键完整非空；每个外键组合都能找到目标。Enrolment.sid 重复与否不决定它是否满足参照完整性；关键是它是否引用有效学生。

\textbf{下文每个操作都从这个初始状态独立判断，不串联执行。} 以下是根据规则进行的静态推演，未在数据库中实际运行。

\subsection{9. 插入：检查新元组与被引用目标}\label{03_relational_model_and_integrity_zh-ux63d2ux5165ux68c0ux67e5ux65b0ux5143ux7ec4ux4e0eux88abux5f15ux7528ux76eeux6807}

来源：第 21--28 页；以下操作使用第 8 节自编数据。

\begin{longtable}[]{@{}
  >{\raggedright\arraybackslash}p{(\linewidth - 2\tabcolsep) * \real{0.5000}}
  >{\raggedright\arraybackslash}p{(\linewidth - 2\tabcolsep) * \real{0.5000}}@{}}
\toprule\noalign{}
\begin{minipage}[b]{\linewidth}\raggedright
拟执行操作
\end{minipage} & \begin{minipage}[b]{\linewidth}\raggedright
判定与原因
\end{minipage} \\
\midrule\noalign{}
\endhead
\bottomrule\noalign{}
\endlastfoot
向 Student 插入 \texttt{(\allowbreak{}S01,\allowbreak{}\ Ada,\allowbreak{}\ ada@example.\allowbreak{}test)\allowbreak{}} & 不合法：sid 已存在，违反键约束 \\
向 Student 插入 \texttt{(\allowbreak{}S03,\allowbreak{}\ Chen,\allowbreak{}\ amy@example.\allowbreak{}test)\allowbreak{}} & 不合法：email 候选键重复，不能只检查主键 \\
向 Enrolment 插入 \texttt{(\allowbreak{}S03,\allowbreak{}\ CS,\allowbreak{}\ 101,\allowbreak{}\ 2026F,\allowbreak{}\ A)\allowbreak{}} & 不合法：Student 没有 S03，违反参照完整性 \\
向 Enrolment 插入 \texttt{(\allowbreak{}S01,\allowbreak{}\ CS,\allowbreak{}\ 201,\allowbreak{}\ 2026F,\allowbreak{}\ B)\allowbreak{}} & 不合法：Course 没有 \texttt{(\allowbreak{}CS,\allowbreak{}201)\allowbreak{}}；不能拆开匹配 \\
向 Enrolment 插入 \texttt{(\allowbreak{}S01,\allowbreak{}\ CS,\allowbreak{}\ 101,\allowbreak{}\ NULL,\allowbreak{}\ B)\allowbreak{}} & 不合法：主键组成属性 term 为 NULL，违反实体完整性 \\
向 Enrolment 插入 \texttt{(\allowbreak{}S01,\allowbreak{}\ CS,\allowbreak{}\ 101,\allowbreak{}\ 2026F,\allowbreak{}\ B)\allowbreak{}} & 不合法：四属性主键已存在，即使 grade 不同也不能再插一行 \\
向 Enrolment 插入 \texttt{(\allowbreak{}S01,\allowbreak{}\ CS,\allowbreak{}\ 101,\allowbreak{}\ 2027S,\allowbreak{}\ NULL)\allowbreak{}} & 合法：term 不同使主键不同，两组外键均有目标，grade 允许 NULL \\
\end{longtable}

答题顺序可以固定为：\textbf{域与必填要求 → 所有候选键 → 主键非空 → 每组外键}。题目要求列出全部违规时，需要查完所有相关规则。

对照课件原例：第 24 页向 RESEARCHER 插入已存在的 Person\#=2，说明键冲突；第 26 页向 COURSE 插入 \texttt{(\allowbreak{}Comp.\allowbreak{}Sci.\allowbreak{},\allowbreak{}\ NULL)\allowbreak{}}，说明复合主键不能缺一部分；第 28 页用不存在的 Supervisee=6 插入 ENROLMENT，说明引用目标缺失。

\subsection{10. 删除：沿外键反向查找引用者}\label{03_relational_model_and_integrity_zh-ux5220ux9664ux6cbfux5916ux952eux53cdux5411ux67e5ux627eux5f15ux7528ux8005}

来源：第 29--30 页。

判断删除时，不只看被删元组自身是否合法，还要看\textbf{有没有其他元组引用它的键}。

\begin{longtable}[]{@{}
  >{\raggedright\arraybackslash}p{(\linewidth - 2\tabcolsep) * \real{0.5000}}
  >{\raggedright\arraybackslash}p{(\linewidth - 2\tabcolsep) * \real{0.5000}}@{}}
\toprule\noalign{}
\begin{minipage}[b]{\linewidth}\raggedright
从初始状态拟删除什么？
\end{minipage} & \begin{minipage}[b]{\linewidth}\raggedright
直接删除的影响
\end{minipage} \\
\midrule\noalign{}
\endhead
\bottomrule\noalign{}
\endlastfoot
Student 中 S01 & 有 Enrolment 引用 S01；若不处理这些引用，违反参照完整性 \\
Course 中 \texttt{(\allowbreak{}CS,\allowbreak{}101)\allowbreak{}} & 两条 Enrolment 引用它；若不处理引用，违反参照完整性 \\
Course 中 \texttt{(\allowbreak{}Math,\allowbreak{}201)\allowbreak{}} & 本例没有记录引用它，删除后仍满足已列约束 \\
Enrolment 中 S01 的那条记录 & 本例没有别的关系引用 Enrolment，删除后仍满足已列约束 \\
\end{longtable}

因此``主键不能被删除''不正确；真正要查的是依赖关系和已声明的处理规则。

第 30 页讨论的处理思路包括：

\begin{itemize}
\tightlist
\item
  \textbf{拒绝删除}：保留被引用元组，避免产生悬空引用。
\item
  \textbf{级联删除引用者}：删除相关元组，但还要检查它们是否继续被其他关系引用。
\item
  \textbf{将引用设为 NULL}：只在外键允许为空且不违反其他约束时可行。
\item
  \textbf{改成另一个有效引用}：新值必须符合域及全部相关约束，不能随意写一个占位编号。
\end{itemize}

在本例中，Enrolment 的两个外键属性组都包含在主键中，因此\textbf{不能用设为 NULL 处理它们}。级联或改引用也必须符合业务含义；不能默认数据库自动选择一种方案。

\subsection{11. 修改：同时考虑旧值与新值}\label{03_relational_model_and_integrity_zh-ux4feeux6539ux540cux65f6ux8003ux8651ux65e7ux503cux4e0eux65b0ux503c}

来源：第 31 页；最后一段为对课件适用范围的解释。

\begin{longtable}[]{@{}
  >{\raggedright\arraybackslash}p{(\linewidth - 4\tabcolsep) * \real{0.3333}}
  >{\raggedright\arraybackslash}p{(\linewidth - 4\tabcolsep) * \real{0.3333}}
  >{\raggedright\arraybackslash}p{(\linewidth - 4\tabcolsep) * \real{0.3333}}@{}}
\toprule\noalign{}
\begin{minipage}[b]{\linewidth}\raggedright
修改位置
\end{minipage} & \begin{minipage}[b]{\linewidth}\raggedright
需要检查什么？
\end{minipage} & \begin{minipage}[b]{\linewidth}\raggedright
本例
\end{minipage} \\
\midrule\noalign{}
\endhead
\bottomrule\noalign{}
\endlastfoot
主键或其他候选键 & 新值是否重复；主键是否变空；旧键是否仍被引用 & 把 Student S01 的 sid 改为 S02 会键冲突；改为 S03 又会让原 S01 引用失去目标，除非同时妥善处理 \\
外键 & 新的完整组合是否能匹配目标；是否允许 NULL；是否兼任主键 & 把 S01 那条选课的 courseNo 从 101 改为 201，会形成不存在的 \texttt{(\allowbreak{}CS,\allowbreak{}201)\allowbreak{}} \\
非键属性 & 仍需检查域、必填和其他业务约束 & 把 grade 改为 Z，虽然没有改键，仍违反本例 grade 允许值规则 \\
\end{longtable}

若将 S01 那条选课的 \texttt{(\allowbreak{}dept,\allowbreak{}courseNo)\allowbreak{}} 一起改为 \texttt{(\allowbreak{}Math,\allowbreak{}201)\allowbreak{}}，新的目标存在，且结果主键不与其他行重复，因此在本例约束下合法。\textbf{改变复合外键时，要检查变更后的整个组合。}

课件把修改主键类比为``删除旧键再插入新键''，便于记住要查旧引用和新值冲突；这是约束分析方法，不要求真的先删除学生再重新插入。

\textbf{课件措辞辨析}：第 31 页 ``Not unless you're modifying the value of a key'' 与 ``all relational integrity constraints are to do with the key values'' 应结合该页讨论的键相关约束理解。它们不能推广成``所有非键更新都安全''：域、必填要求及业务规则同样可能被非键更新破坏。

\subsection{12. 易混点速查}\label{03_relational_model_and_integrity_zh-ux6613ux6df7ux70b9ux901fux67e5}

\begin{longtable}[]{@{}
  >{\raggedright\arraybackslash}p{(\linewidth - 2\tabcolsep) * \real{0.5000}}
  >{\raggedright\arraybackslash}p{(\linewidth - 2\tabcolsep) * \real{0.5000}}@{}}
\toprule\noalign{}
\begin{minipage}[b]{\linewidth}\raggedright
容易混淆
\end{minipage} & \begin{minipage}[b]{\linewidth}\raggedright
正确区分
\end{minipage} \\
\midrule\noalign{}
\endhead
\bottomrule\noalign{}
\endlastfoot
Relation / ER relationship & 元组集合 / 实体之间的联系 \\
Schema / instance & 结构和规则 / 某时刻的内容 \\
Degree / cardinality & 属性数 / 当前元组数 \\
关系的 cardinality / ER cardinality ratio & 实际行数 / 关联数量的业务上限 \\
无完全重复元组 / 键唯一 & 整行不重复 / 指定属性组不能在不同行中重复 \\
Candidate key / primary key & 所有最小超键之一 / 被选定的一个候选键 \\
主键 / 外键 & 标识本关系元组 / 引用目标关系的键；同一属性可兼任 \\
外键值重复 / 悬空引用 & 多行引用同一目标通常允许 / 引用不存在的目标不允许 \\
复合外键匹配 / 分别查每列 & 同一目标元组中整体匹配 / 分别存在不足以证明合法 \\
Key constraint / entity integrity & 唯一性 / 主键非空 \\
NULL / 0 / 空字符串 & 缺失标记 / 一个数字值 / 一个字符串值 \\
约束允许的状态 / 业务完全正确 & 满足已声明规则 / 还取决于是否把需要的业务规则声明完整 \\
\end{longtable}

\subsection{13. 复习清单与英文表达}\label{03_relational_model_and_integrity_zh-ux590dux4e60ux6e05ux5355ux4e0eux82f1ux6587ux8868ux8fbe}

以下是自我复习目标，不是正式模拟测试或官方分值安排。

\begin{itemize}
\tightlist
\item[$\square$]
  从一组表中指出 relation、tuple、attribute、domain、schema 和 instance。
\item[$\square$]
  给出关系的 degree 与 cardinality，并解释空实例。
\item[$\square$]
  说明交换显示列与交换元组值的区别。
\item[$\square$]
  区分无重复行、超键、候选键和主键。
\item[$\square$]
  解释为何不能无条件使用 \texttt{2\^{}n\ −\ 1} 计算超键数。
\item[$\square$]
  标出每组外键的引用方向，整体核对复合外键。
\item[$\square$]
  分别解释键约束、实体完整性、参照完整性。
\item[$\square$]
  根据课件规则判断外键全空、部分空及非空不匹配情况。
\item[$\square$]
  从同一初始状态独立判断插入、删除与修改，说明具体冲突元组或缺失目标。
\item[$\square$]
  说明级联删除与设为 NULL 的前提，不把它们当作自动行为。
\end{itemize}

可用于英文简答的表达：

\begin{quote}
A relation is a set of tuples defined over a collection of named attributes.
\end{quote}

\begin{quote}
A relation schema describes the structure; a relation instance contains the tuples at a particular time.
\end{quote}

\begin{quote}
A candidate key is a minimal superkey, and a primary key is a selected candidate key.
\end{quote}

\begin{quote}
Entity integrity requires every component of a primary key to be non-NULL.
\end{quote}

\begin{quote}
Referential integrity requires a non-NULL foreign-key value to match a referenced key value.
\end{quote}

\begin{quote}
A composite foreign key must match the corresponding key components in the same referenced tuple.
\end{quote}

\begin{quote}
Deleting a referenced tuple requires checking the tuples that refer to it.
\end{quote}

\subsection{来源与整理记录}\label{03_relational_model_and_integrity_zh-ux6765ux6e90ux4e0eux6574ux7406ux8bb0ux5f55}

\begin{itemize}
\tightlist
\item
  教师主课件：\href{../../../slides/Lecture2A.pptx}{Lecture2A.pptx}，主要第 4--32 页。已读取文字并查看第 7、10、15、17、19、24、26、28--29、31 页原图，核对表头下划线、复合外键与操作例子；第 33--55 页的映射步骤不在本篇展开。
\item
  补充对照：\href{../../../slides/SQLpart1.pptx}{SQLpart1.pptx}，第 31--32 页，仅用于区分经典集合语义与 SQL 重复值；不在此教授 SQL 查询。
\item
  范围依据：\href{../../../syllabus/CSC5041_Teaching_Plan.pdf}{CSC5041\_Teaching\_Plan.pdf}，第 1 页 Week 3。degree/cardinality、域约束与业务约束的解释、自编模式与数据、课件措辞的限定均已与原文区分。
\item
  本版为 2026-10-20 的 LaTeX 重排；示例证据状态沿用原说明，未新增数据库实测。
\item
  输入版本 SHA-256：Lecture2A \texttt{\seqsplit{cecf28254342cf35e9a0b1b86a007fba728ed17a900e9e05182f2f2616081ee0}}；SQLpart1 \texttt{\seqsplit{aa28c989fb3674e780418c78c483bba20ff5a50b7da067b0341e60410a68c8ed}}；syllabus 教学计划 \texttt{\seqsplit{31380a7fe6504e46541b6db0bb2f24859c69188bfeb9681dca73354ebeb68295}}。
\end{itemize}
